\section{Molecular and constitutional mechanisms}

\subsection{CACNA1B and the developmental role of CaV2.2}

CACNA1B encodes the pore-forming component of CaV2.2, an N-type calcium channel implicated in presynaptic release. CaV2.1, encoded by CACNA1A, is another major release channel. Channel contributions vary across neurons, compartments, developmental stages, and preparations.

Rat auditory-brainstem experiments demonstrate developmental changes in the channel types supporting transmission \cite{ref8}. These data justify a synapse-specific handoff visualization, not a universal schedule in which an immature brain is CaV2.2 and a mature brain is CaV2.1. A different preparation may retain substantial N-type contribution; absence of the illustrated switch need not mean arrested development.

Human biallelic CACNA1B loss-of-function variants were reported in a severe epilepsy–dyskinesia syndrome \cite{ref9}. This demonstrates that major disruption can matter profoundly. It does not show that ordinary birth variation produces a smaller version of the same syndrome. The proposed environmental bridge needs direct evidence in a biologically plausible exposure range. Hypoxia-specific CACNA1B splice claims and a birth-trauma-induced universal stalled switch are not asserted.


\subsection{The \ensuremath{\alpha_2\delta} / CACNA2D developmental hinge}

The \ensuremath{\alpha_2\delta} family supplies a particularly informative interface because trafficking and synapse organization can be experimentally separated. Eroglu et al. identified \ensuremath{\alpha_2\delta}-1 as a neuronal thrombospondin receptor supporting excitatory synaptogenesis \cite{ref10}. Risher et al. investigated its postsynaptic Rac1-dependent synaptogenic role in developing mouse cortex \cite{ref11}. Ferron et al. experimentally connected proteolytic maturation of \ensuremath{\alpha_2\delta} with vesicular release probability \cite{ref12}.

These findings support multiple molecular functions, not a single scalar “\ensuremath{\alpha_2\delta} level.” CACNA2D1–4 encode distinct family members; results for \ensuremath{\alpha_2\delta}-1 cannot be transferred automatically to every isoform. Synaptogenic and channel-related functions can overlap without being identical. Astrocyte-dependent synapse formation should not be silently reduced to presynaptic calcium influx.

A mouse nerve-injury study found \ensuremath{\alpha_2\delta}-1-dependent CaV2.2 trafficking in selected dorsal-root neurons \cite{ref29}; that preparation is not a perinatal brain model. Mouse brain expression profiles vary by region and developmental stage, and one activity-blockade condition did not alter the measured channel-expression program \cite{ref30}. A biochemical assay found that gabapentin did not directly inhibit binding of soluble thrombospondin-4 to \ensuremath{\alpha_2\delta}-1 \cite{ref31}, even though earlier synaptogenesis experiments found drug-sensitive outcomes \cite{ref10}. The binding event, channel trafficking and downstream synapse formation therefore remain separate steps. Rare biallelic CACNA2D1 loss-of-function in two humans impaired CaV2 trafficking in assays and was associated with severe early encephalopathy \cite{ref35}; it cannot identify the effects of common variation or ordinary exposure.

The unproven bridge is explicit: do realistic endocrine, inflammatory, metabolic, or microbial perturbations alter particular \ensuremath{\alpha_2\delta} states during relevant sensitive periods, and do those changes mediate durable network differences? Hinge status would require temporal ordering, intervention, selective disruption or rescue, and a comparison with alternative pathways. A correlation between exposure and CACNA2D transcript abundance alone is insufficient.


\subsection{The RCCX constitutional interface}

The broad project hypothesis invokes immune, endocrine, and connective-tissue convergence at the RCCX region \cite{projP1}. Specific human evidence exists for contiguous CYP21A2–TNXB disruption and the CAH-X phenotype: Merke et al. documented connective-tissue manifestations associated with tenascin-X haploinsufficiency in a subset of congenital adrenal hyperplasia patients \cite{ref13}. This is a defined molecular-clinical relationship.

Additional clinical genetics work maps the RCCX structural arrangement and CAH-X testing boundary \cite{ref22}. Complement C4 variation has a distinct schizophrenia association and a mouse postnatal-pruning result \cite{ref23}, not a general RCCX developmental pathway. ADHD genome-wide analysis identified 27 risk loci \cite{ref24}, directly challenging a one-locus account. In a selected 102-person hEDS/HSD clinical cohort, 48\% met a POTS definition \cite{ref33}; this is not a population prevalence or evidence that RCCX caused the autonomic phenotype.

It must not be generalized into an established RCCX cause of autism, ADHD, overexcitability, autonomic illness, or polymathic capacity. Chromosomal proximity does not demonstrate coordinated expression, a common latent syndrome, or causal mediation through \ensuremath{\alpha_2\delta}. No direct RCCX \ensuremath{\rightarrow} \ensuremath{\alpha_2\delta} mechanism is established by the evidence used here.

The research design should compare the incremental value of prespecified RCCX variables against broader polygenic, obstetric, immune, endocrine, and connective-tissue measures. If the RCCX subset adds no replicated explanatory value, the RCCX-specific module should be rejected while the broader constitution \ensuremath{\times} development question remains testable. Defining susceptibility retrospectively from the very outcome to be explained would make the model circular.
